\subsection{主理想整环中的整除性}

设 A 为主理想整环，并设 K 为其分式域（I，第 116 页）；我们将看到由 K 的主分式理想（VI，第 6 页）组成的序群 \(P^{*}\) 是格序的；更精确地说：

命题1. --- 设 K 为主理想整环 A 的分式域，且设 \((x_{\iota})_{\iota \in I}\) 是 K 中具有公分母 \(b \in K^{*}\) 的元素族（换句话说， \(bx_{\iota} \in A\) 对所有 \(\iota\) ). 那么：

\begin{enumerate}
\def\labelenumi{\alph{enumi})}
\item
  族 \((x_{\iota})\) 在 K 中有一个最大公因子。
\item
  \((x_{i})\) 的每个最大公因子都可以表示为 \(d = \sum_{i}a_{i}x_{i}\) ，其中
\end{enumerate}

\(a_{i}\) 是A的元素，其中除有限多个外均为零。

事实上，理想 \(\sum_{i} A b x_{i}\) （ \(A\) 的）是主理想，因此具有形式 \(A d'\) 。令

\(d' = bd (d \in \mathbf{K})\) 。由关系式 \(d' = \sum_{i} a_i b x_i\) ，我们推导出 \(d = \sum_{i} a_i x_i\) ，其中

\(a_{i} \in A\) 。因此， \(x_{i}\) 的每个公因子都整除 d。另一方面，由于按构造 bd 是 \(bx_{i}\) 的公因子，由此可知 d 是 \(x_{i}\) 的公因子.

注：--- 命题 1 无限制地适用于 A 的元素的任意族 \((x_{i})\) （取 \(b = 1\) ），也适用于 K 中元素的任意有限族 \((x_{i})\) （若 \(x_{i} = c_{i}b_{i}^{-1}\) ，其中 \(c_{i} \in \mathbf{A}\) ， \(b_{i} \in \mathbf{A}\) ，则取 \(b\) 为诸 \(b_{i}\) 的乘积）.

推论。——设 \((x_{t})\) 是主理想整环 A 中元素的一个任意族，其中 A 作为子环包含在整环 B 中，并设 d 为族 \((x_{t})\) 在 A 中的一个最大公因子。则族 \((x_{t})\) 在 B 中有最大公因子，且 d 是其中之一。

确实， \(d\) 是 \(x_{i}\) 在 B 中的公因子。另一方面，关系式 \(d = \sum_{i} a_{i} x_{i}\) 表明 \(x_{i}\) 在 B 中的每个公因子都整除 \(d\) .

这个推论的一个重要应用是当 \(\mathbf{A} = \mathbf{K}[\mathbf{X}]\) 且 \(\mathbf{B} = \mathbf{E}[\mathbf{X}]\) 时，其中 \(\mathbf{K}\) 是一个域，且 \(\mathbf{E}\) 是 \(\mathbf{K}\) 的一个扩张（IV，第 12 页，推论 1）。

命题1的第一个断言表明序群 \(P^{*}\) 是格序的（VI，第 10 页）。特别地，K 的元素的每个有限族都有最小公倍元。因此我们可以将VI第10至17页中记为（DIV）的结果应用于主理想整环。

以下结果是命题1第二个断言的推论：

定理1（“贝祖恒等式”）。--- 主理想整环 A 中的元素 \(x_{\iota}\) ( \(\iota \in I\) ）是整体互素的，当且仅当存在 A 的元素 \(a_{\iota}\) ( \(\iota \in I\) ）（除有限多个外均为零），使得 \(\sum a_{\iota}x_{\iota} = 1\) .

该条件的必要性正是命题 1。反之，如果 \(\sum_{i} a_i x_i = 1\) ，那么 \(x_i\) 的每个公因子都整除 1，因此 1 是 \(x_i\) 的一个最大公因子.

命题 2. --- 设 \(a, b, d, m\) 与 \(p\) 是分式域 \(K\) （主理想整环 \(A\) 的）的元素.

\begin{enumerate}
\def\labelenumi{\alph{enumi})}
\item
  “d 是 a 和 b 的一个最大公因子”等价于“(d) = (a) + (b)”。
\item
  “m 是 a 和 b 的一个最小公倍元”等价于“m = (a) ∩ (b)”。
\item
  «p是A的不可约元素»等价于«(p)是A的非零极大理想»，也等价于«(p)是A的非零素理想»。
\end{enumerate}

我们已经证明了 \(a\) )（命题 1）。由于 \(a\) 与 \(b\) 的公倍数就是 \((a) \cap (b)\) 的元素，并且由于根据假设，理想 \((a) \cap (b)\) 是主理想，设 \((a) \cap (b) = (m)\) ，由此可得 \(m\) 是 \(a\) 与 \(b\) 的一个最小公倍元，这证明了 \(b\) )。最后，说 \(p \neq 0\) 是 \(A\) 的不可约元素，按定义（VI，第 17 页）意味着 \((p)\) 是 \(\neq A\) 的主理想（ \(A\) 中的）按包含关系排序所成族的一个极大元；由于 \(A\) 的理想都是主理想，这意味着 \((p)\) 是 \(A\) 的一个极大理想，从而得到 \(c\) )，根据 VI 第 17 页的注记。

在主理想整环 A 中，有限个理想的加法（相应地，交）也称为这些理想的 gcd（相应地，lcm）。

命题 3. --- 设 \(a, b, c\) 是主理想整环 A 的分式域的元素，并令 \(d\) 为 \(a\) 与 \(c\) 的一个最大公因子；则同余式 \(ax \equiv b \pmod{c}\) 有解 \(x_0 \in A\) 当且仅当 \(d\) 整除 \(b\) ；在这种情况下，A 的元素 \(x\) （满足 \(ax \equiv b \pmod{c}\) 者）恰好是那些满足 \(x \equiv x_0 \pmod{cd^{-1}}\) 的元素.

如果 \(ax \equiv b \pmod{c}\) ，其中 \(x \in \mathbf{A}\) ，则存在 \(y \in \mathbf{A}\) 使得 \(b = ax + cy\) ，所以 \(d\) 整除 \(b\) 。反之，若 \(d\) 整除 \(b\) ，则 \(b = ax_0 + cy_0\) ，其中 \(x_0, y_0 \in \mathbf{A}\) （命题 1），因此 \(ax_0 \equiv b \pmod{c}\) ；此外，关系 \(ax \equiv b \pmod{c}\) 便等价于 \(a(x - x_0) \equiv 0 \pmod{c}\) ；令 \(a = da'\) ， \(c = dc'\) ，后者变为 \(a'(x - x_0) \equiv 0 \pmod{c'}\) 。但这（对于 \(x \in \mathbf{A}\) ）等价于 \(x - x_0 \equiv 0 \pmod{c'}\) ，因为 \(a'\) 与 \(c'\) 互素（VI，第14页，命题10（DIV）及VI，第15页，命题11的推论2（DIV））。

命题4. --- 设 \((a_i)_{1 \leq i \leq n}\) 是主理想整环 A 中两两互素的元素组成的有限族。则从 \(\mathbf{A}/\left(\prod_{i=1}^{n} a_i\right)\) 到乘积 \(\prod_{i=1}^{n} \mathbf{A}/(a_i)\) 的典范同态（I，第 110 页）是 A-代数的同构。

这由I，第110页，命题9以及命题2 a)得出。

命题4的结论在未假设 \(a_{i}\) 两两互素的情况下不成立（参见VII，第24页，命题9）。

